\documentclass[11pt,openany]{book}
\usepackage[T1]{fontenc}
\usepackage[margin=1in]{geometry}
\usepackage{makeidx}
\usepackage{hyperref}
\hypersetup{colorlinks=true,linkcolor=blue,citecolor=blue,urlcolor=blue}
\makeindex

\title{A Pocket Glossary of Typesetting}
\author{R. Okafor}
\date{}

\begin{document}
\maketitle

\chapter{Boxes, glue, and penalties}
\label{ch:main}

\section{Boxes and glue}

Every page is built from \index{box}boxes joined by
\index{glue}glue. A horizontal box\index{box!horizontal} holds a line of
type, while a vertical box\index{box!vertical} stacks lines into a page.
Glue\index{glue!stretch} stretches to fill short lines, and
glue\index{glue!shrink} shrinks to fit long ones. The penalties of
Section~\ref{sec:breaks} on page~\pageref{sec:breaks} decide where these
lines may end.

\section{Penalties and breaks}
\label{sec:breaks}

\index{penalty}Penalties mark good and bad places to break. A forced break
is just an infinite\index{penalty!infinite} penalty elsewhere, and the
\index{badness}badness of a line measures how far the glue had to give.
Widows\index{widow} and orphans\index{orphan} are controlled the same way,
through penalties\index{penalty!widow} on the page builder.

\printindex
\end{document}
